\documentclass[9pt,a4paper,twocolumn]{ctexart}

% --- Packages ---
\usepackage[top=1.2cm, bottom=1.2cm, left=1.2cm, right=1.2cm, columnsep=0.8cm]{geometry}
\usepackage{hyperref}
\usepackage{graphicx}
\usepackage{float}
\usepackage{longtable}
\usepackage{tabularx}
\usepackage{booktabs}
\usepackage{enumitem}
\usepackage{xcolor}
\usepackage{fancyhdr}
\usepackage{listings}
\usepackage{fontspec}
\usepackage{titlesec}
\usepackage{caption}

% --- Code block styling ---
\definecolor{codebg}{RGB}{245,245,245}
\definecolor{codeframe}{RGB}{220,220,220}
\definecolor{codekw}{RGB}{0,0,192}
\definecolor{codecomment}{RGB}{0,128,0}
\definecolor{codestring}{RGB}{163,21,21}
\definecolor{codenum}{RGB}{128,0,128}
\definecolor{codepunct}{RGB}{90,90,90}
\definecolor{badgegreen}{RGB}{34,139,34}
\definecolor{badgeblue}{RGB}{30,144,255}
\definecolor{badgeblack}{RGB}{50,50,50}

% --- Difficulty badges (rounded corners via tikz, pre-rendered in saveboxes) ---
\usepackage{tikz}
\newsavebox{\badgechubox}
\savebox{\badgechubox}{\tikz[baseline]{\node[fill=badgegreen,text=white,rounded corners=2.5pt,inner xsep=4pt,inner ysep=1.5pt,font=\sffamily\footnotesize\bfseries]{初级};}}
\newsavebox{\badgezhongbox}
\savebox{\badgezhongbox}{\tikz[baseline]{\node[fill=badgeblue,text=white,rounded corners=2.5pt,inner xsep=4pt,inner ysep=1.5pt,font=\sffamily\footnotesize\bfseries]{中级};}}
\newsavebox{\badgegaobox}
\savebox{\badgegaobox}{\tikz[baseline]{\node[fill=badgeblack,text=white,rounded corners=2.5pt,inner xsep=4pt,inner ysep=1.5pt,font=\sffamily\footnotesize\bfseries]{高级};}}
\newcommand{\lvchu}{\usebox{\badgechubox}}
\newcommand{\lvzhong}{\usebox{\badgezhongbox}}
\newcommand{\lvgao}{\usebox{\badgegaobox}}
\pdfstringdefDisableCommands{%
  \def\lvchu{[初级]}%
  \def\lvzhong{[中级]}%
  \def\lvgao{[高级]}%
}

\lstset{
  basicstyle=\ttfamily\scriptsize,
  keywordstyle=\color{codekw}\bfseries,
  commentstyle=\color{codecomment}\itshape,
  stringstyle=\color{codestring},
  backgroundcolor=\color{codebg},
  frame=single,
  rulecolor=\color{codeframe},
  breaklines=true,
  breakatwhitespace=false,
  breakindent=0pt,
  postbreak=\mbox{\textcolor{red}{$\hookrightarrow$}\space},
  numbers=left,
  numberstyle=\tiny\color{gray},
  columns=flexible,
  keepspaces=true,
  showstringspaces=false,
  tabsize=2,
  aboveskip=4pt,
  belowskip=4pt,
  extendedchars=true,
  literate={，}{{，}}1 {；}{{；}}1 {！}{{！}}1 {？}{{？}}1 {“}{{“}}1 {”}{{”}}1 {（}{{（}}1 {）}{{）}}1 {《}{{《}}1 {》}{{》}}1,
}

% --- Extra language definitions (no built-in listings support) ---
\lstdefinelanguage{json}{
  morekeywords={true,false,null},
  sensitive=true,
  morestring=[b]",
  comment=[l]{//},
  morecomment=[s]{/*}{*/},
}
\lstdefinelanguage{yaml}{
  morekeywords={true,false,null,yes,no,on,off},
  sensitive=false,
  morecomment=[l]{\#},
  morestring=[b]",
  morestring=[d]',
}
\lstdefinelanguage{http}{
  morekeywords={GET,POST,PUT,DELETE,PATCH,HEAD,OPTIONS,HTTP,Host,Accept,Authorization,Content-Type,Content-Length,User-Agent},
  sensitive=true,
  morecomment=[l]{\#},
  morestring=[b]",
}
\lstdefinelanguage{TypeScript}{
  morekeywords={abstract,any,as,async,await,break,case,catch,class,const,continue,debugger,declare,default,delete,do,else,enum,export,extends,false,finally,for,from,function,get,if,implements,import,in,instanceof,interface,keyof,let,module,namespace,never,new,null,number,of,package,private,protected,public,readonly,return,set,static,string,super,switch,symbol,this,throw,true,try,type,typeof,undefined,unknown,var,void,while,with,yield,Record,Array,Promise,AsyncIterable,ReadableStream,Date,Error,Object,JSON},
  sensitive=true,
  morecomment=[l]{//},
  morecomment=[s]{/*}{*/},
  morestring=[b]",
  morestring=[b]',
  morestring=[b]`{`},
}

% --- Compact spacing ---
\linespread{0.95}
\setlength{\parskip}{1pt plus 1pt minus 1pt}
\setlength{\parindent}{0pt}

% --- Table styling ---
\renewcommand{\arraystretch}{1.1}

% --- Heading styling (numbered) ---
\titleformat{\section}{\normalsize\bfseries}{\thesection\quad}{0em}{}
\titlespacing{\section}{0pt}{6pt}{2pt}
\titleformat{\subsection}{\small\bfseries}{\thesubsection\quad}{0em}{}
\titlespacing{\subsection}{0pt}{4pt}{1pt}
\titleformat{\subsubsection}{\small\bfseries}{\thesubsubsection\quad}{0em}{}
\titlespacing{\subsubsection}{0pt}{3pt}{1pt}

% --- Page style ---
\pagestyle{fancy}
\fancyhf{}
\fancyhead[L]{\small\emph{\leftmark}}
\fancyhead[R]{\small\thepage}
\renewcommand{\headrulewidth}{0.4pt}
\renewcommand{\sectionmark}[1]{}  % prevent sections from overwriting leftmark

% --- Hyperref setup ---
\hypersetup{
  colorlinks=true,
  linkcolor=blue,
  urlcolor=blue,
  citecolor=blue,
}

% --- Caption format ---
\DeclareCaptionFormat{myformat}{\small\bfseries #1#2#3}
\captionsetup{format=myformat}

\begin{document}

\title{AI 应用系统设计：从 Prompt Demo 到生产级架构}
\date{}
\maketitle
\markboth{Java 版 Agent 知识}{ Java 版 Agent 知识}

\section*{Demo 架构为什么扛不住生产流量}


先看一个最常见的 Demo：



\begin{lstlisting}
前端输入问题 -> 后端拼 Prompt -> 调用模型 API -> 返回答案
\end{lstlisting}



这条链路能演示产品想法，但它缺了生产系统最关键的 6 件事。


{\footnotesize
\begin{tabular}{|p{\dimexpr 0.075\columnwidth-2\tabcolsep\relax}|p{\dimexpr 0.245\columnwidth-2\tabcolsep\relax}|p{\dimexpr 0.679\columnwidth-2\tabcolsep\relax}|}
\hline
\textbf{维度} & \textbf{Prompt Demo} & \textbf{生产级架构} \\
\hline
稳定性 & 单模型、单调用，失败就报错 & 多模型路由、重试、fallback、熔断、降级响应 \\
\hline
权限 & 默认用户能问什么就查什么 & 检索前权限过滤，工具调用按用户和租户鉴权 \\
\hline
成本 & 只看一次调用能不能成功 & Token 预算、模型分层、缓存、成本归因和限额 \\
\hline
可观测 & 记录用户问题和最终答案 & 记录 Prompt、检索片段、工具调用、模型输出、Token、延迟、错误 \\
\hline
评测 & 靠人工试几条样例 & 固定评测集、线上抽样、LLM-as-Judge、人工复核闭环 \\
\hline
数据治理 & 文档直接入库，日志随便存 & PII 脱敏、数据留存、审计、版本化、删除和授权链路 \\
\hline
\end{tabular}
}



AI 应用把部分决策交给概率模型，故障不一定落在某一行代码。LLM 出错可能来自 Prompt 版本、上下文顺序、检索噪声、工具描述、模型采样、权限过滤或输出解析中的任何一环。生产级架构必须把输入、执行、输出和反馈全部工程化，保证可追踪、可回放、可评测、可治理。



补充阅读：大模型 API 调用工程实践（流式输出、重试、限流与结构化返回）；Token、上下文窗口和采样参数基础见 LLM 运行机制。



\section*{生产级 AI 应用的标准分层架构}


推荐按职责拆层。不同公司命名会有差异，但生产系统里的边界大体一致。



\subsection*{入口层：把用户请求变成可治理的任务}


入口层不能只当 Controller 用。它至少要做这些事：


\begin{itemize}[itemsep=1pt, topsep=2pt]
  \item 认证鉴权：确认用户、租户、角色、数据范围。
  \item 请求标准化：把 Web、App、API、Webhook、定时任务统一成内部任务模型。
  \item 限流与防刷：按用户、租户、模型能力和业务场景限流。
  \item 幂等控制：异步任务、工具调用、支付类操作必须有幂等键。
  \item 敏感内容预处理：PII 脱敏、恶意输入检测、Prompt 注入初筛。
\end{itemize}


入口层最后应该产出结构化请求，而不是只把用户输入当成一段字符串往后传：



\begin{lstlisting}[language=Java]
public record AiRequest(
        String requestId,
        String tenantId,
        String userId,
        String sceneCode,
        String input,
        Map<String, Object> variables,
        PermissionScope permissionScope
) {
}
\end{lstlisting}



\subsection*{业务编排层：决定这次请求怎么跑}


业务编排层负责决定这次请求怎么执行：


\begin{itemize}[itemsep=1pt, topsep=2pt]
  \item 这次是普通问答、RAG 问答、Agent 多步任务，还是批处理任务？
  \item 需要哪些上下文：历史会话、用户画像、知识库、实时业务数据？
  \item 是否允许调用工具？哪些工具需要二次确认？
  \item 应该走同步、流式，还是异步？
  \item 输出要不要进入评测、人工审核或后处理？
\end{itemize}


这层别把所有逻辑都塞进一个“超级 Prompt”。能确定的规则用代码处理，无法穷举的语言理解再交给模型。边界清楚，系统才容易排查。



\subsection*{模型网关：把模型调用变成基础设施}


模型网关负责统一接入 OpenAI、Anthropic、Google Gemini、私有化模型、Embedding 模型、Rerank 模型等能力。它隐藏不同 API 的差异，对上提供稳定接口。模型网关本身可以单独展开一篇，细节可以看 大模型网关详解：多模型路由、Fallback、限流与成本控制。



\begin{figure}[H]
\centering
\includegraphics[width=0.85\textwidth]{/home/mao/demos/latex_demos/ai-agent-knowledge/HTML/AI 应用系统设计：从 Prompt Demo 到生产级架构 _ JavaGuide_files/llm-gateway-overview.jpg}
\caption*{LLM 网关示意图}
\end{figure}



模型网关的核心能力包括：


\begin{itemize}[itemsep=1pt, topsep=2pt]
  \item 多模型路由：按场景、成本、延迟、语言、上下文长度和成功率选择模型。
  \item fallback：主模型失败、超时、限额不足时切到备用模型。
  \item 限流与熔断：避免供应商异常拖垮业务线程池。
  \item Token 预算：估算输入输出 Token，超预算时压缩上下文或降级模型。
  \item 成本归因：按租户、用户、场景、Prompt 版本记录成本。
  \item 统一观测：记录模型请求、响应、错误、TTFT、总延迟、Token usage。
\end{itemize}


OpenAI、Anthropic、Google 等官方文档都在持续更新模型、工具、流式、评测和成本相关能力。涉及具体模型名、上下文窗口、价格、可用区域和工具支持时，建议在配置中心或模型注册表里维护，并标注“以官方文档最新展示为准”，不要写死在业务代码里。



\subsection*{Prompt 与 Context 管理：不要把 Prompt 当代码里的字符串}


Prompt 在生产环境里应该被当成一种可版本化配置，不能散落成代码里的多行字符串。



它至少需要支持：


\begin{itemize}[itemsep=1pt, topsep=2pt]
  \item 模板版本：每次修改生成新版本，旧版本可回放。
  \item 变量注入：业务变量、用户输入、检索结果、工具结果分区注入。
  \item 灰度发布：按租户、用户比例、场景开关选择 Prompt 版本。
  \item 快速回滚：线上效果变差时能切回稳定版本。
  \item 审计记录：谁在什么时间改了什么，为什么改。
  \item 运行时绑定：每次请求记录使用的 Prompt 名称、版本和变量摘要。
\end{itemize}


一个很实用的规则：\textbf{Prompt 变更要像代码变更一样可追踪，但发布频率可以比代码更高}。



Langfuse 官方文档把 Prompt Management、Tracing、Evaluation 放在同一套 LLM 工程平台中，原因相同：Prompt 不只影响生成文本，还影响检索、工具调用、成本和评测结果。不强制使用 Langfuse，但这类数据应在自己的系统中串起来。



Prompt 写法见 大模型提示词工程（Prompt Engineering）是什么？提示词技巧有哪些？。“哪些信息该进上下文、进多少、什么时候压缩”见 上下文工程（Context Engineering）是什么？和 Prompt Engineering 有什么区别？。



\begin{figure}[H]
\centering
\includegraphics[width=0.85\textwidth]{/home/mao/demos/latex_demos/ai-agent-knowledge/HTML/AI 应用系统设计：从 Prompt Demo 到生产级架构 _ JavaGuide_files/context-engineering-vs-context-engineering-dimension-comparison.jpg}
\caption*{Context Engineering 和 Prompt Engineering 差别}
\end{figure}



\subsection*{RAG、Memory、Tool：三类上下文不要混在一起}


很多 AI 系统越做越乱，是因为把所有信息都叫“上下文”。



推荐按以下边界拆分：


{\footnotesize
\begin{tabular}{|p{\dimexpr 0.109\columnwidth-2\tabcolsep\relax}|p{\dimexpr 0.400\columnwidth-2\tabcolsep\relax}|p{\dimexpr 0.145\columnwidth-2\tabcolsep\relax}|p{\dimexpr 0.345\columnwidth-2\tabcolsep\relax}|}
\hline
\textbf{类型} & \textbf{存什么} & \textbf{生命周期} & \textbf{核心风险} \\
\hline
RAG & 企业文档、产品手册、制度、代码文档、工单知识 & 由知识库更新决定 & 检索不到、越权召回、过期文档、引用错配 \\
\hline
Memory & 用户偏好、历史决策、长期画像、任务经验 & 随用户和会话演化 & 错误记忆固化、隐私泄露、过时记忆干扰 \\
\hline
Tool & 查询订单、创建工单、发邮件、改配置、查数据库 & 运行时按需调用 & 参数错误、权限越界、敏感操作误执行 \\
\hline
\end{tabular}
}



三者底层都可能用向量检索、结构化存储和重排，但服务目标完全不同。RAG 提供共享知识源，Memory 提供个性化背景，Tool 连接真实业务系统。



\begin{figure}[H]
\centering
\includegraphics[width=0.85\textwidth]{/home/mao/demos/latex_demos/ai-agent-knowledge/HTML/AI 应用系统设计：从 Prompt Demo 到生产级架构 _ JavaGuide_files/agent-memory-rag-vs-memory.pdf}
\caption*{长期记忆与 RAG（检索增强生成）的区别}
\end{figure}



\textbf{高频盲区：不要把 Memory 当成个人版 RAG 随便塞。} 记忆一旦写错，后续每轮都会被污染。生产环境里的 Memory 写入通常要异步执行，并经过 Schema 校验、置信度过滤、过期策略和人工审核入口。



RAG 的基础概念可以从 万字详解 RAG 基础概念 看起；文档如何解析、清洗和切 Chunk，可以看 RAG 文档处理与切分策略；检索效果调优看 万字详解 RAG 优化：从召回、重排到上下文工程的系统调优。Memory 单独展开的话，可以看 AI Agent 记忆系统：短期记忆、长期记忆与记忆演化机制。



\section*{同步、流式、异步三种交互模式怎么选}


AI 应用不是所有请求都适合 HTTP 同步等待。交互模式选错，用户体验和系统稳定性都会被拖垮。


{\footnotesize
\begin{tabular}{|p{\dimexpr 0.049\columnwidth-2\tabcolsep\relax}|p{\dimexpr 0.259\columnwidth-2\tabcolsep\relax}|p{\dimexpr 0.173\columnwidth-2\tabcolsep\relax}|p{\dimexpr 0.185\columnwidth-2\tabcolsep\relax}|p{\dimexpr 0.333\columnwidth-2\tabcolsep\relax}|}
\hline
\textbf{模式} & \textbf{适合场景} & \textbf{优势} & \textbf{风险} & \textbf{后端设计要点} \\
\hline
同步请求 & 短问答、分类、抽取、低延迟小任务 & 实现简单，调用链清晰 & 超时敏感，容易占满线程 & 设置短超时、快速失败、结果缓存 \\
\hline
流式响应 & 聊天、长答案、代码生成、语音前置文本 & 首字体验好，用户感知等待更短 & 中途失败处理复杂，前端状态更多 & SSE/WebSocket、TTFT 监控、可取消生成 \\
\hline
异步任务 & 报告生成、批量评测、长文档分析、多工具任务 & 可排队、可重试、可恢复 & 任务状态和通知链路复杂 & 任务表、队列、进度事件、幂等和补偿 \\
\hline
\end{tabular}
}



可以先按这个经验阈值选：


\begin{itemize}[itemsep=1pt, topsep=2pt]
  \item \textbf{能在 3 秒内稳定完成的任务}，优先同步。这个值不是标准答案，要结合网关、负载均衡和客户端超时一起定。
  \item \textbf{用户需要立刻看到模型开始输出的任务}，优先流式。
  \item \textbf{依赖长文档、多轮工具调用或批量处理的任务}，优先异步。
\end{itemize}


别为了“看起来像 ChatGPT”把所有接口都做成流式。比如标签分类、风险评分、路由决策这类内部调用，流式收益不大，反而会增加链路复杂度。



流式输出、重试、限流和结构化返回在 大模型 API 调用工程实践 里有更完整的工程拆解。如果场景是实时语音，还要考虑 VAD、ASR、TTS、打断和端到端延迟，可以继续看 AI 语音技术详解：从 ASR、TTS 到实时语音 Agent 的工程化落地。



\section*{Prompt 管理：从模板字符串到版本系统}


生产级 Prompt 管理可以先按 5 个对象建模：


\begin{itemize}[itemsep=1pt, topsep=2pt]
  \item \texttt{prompt\_template}：Prompt 基本信息，例如名称、场景、类型、状态。
  \item \texttt{prompt\_version}：具体内容、变量定义、模型参数、创建人、变更说明。
  \item \texttt{prompt\_release}：某个版本发布到哪个环境、哪些租户、多少流量。
  \item \texttt{prompt\_run}：每次调用绑定的 Prompt 版本、变量摘要和模型输出。
  \item \texttt{prompt\_eval\_result}：某个 Prompt 版本在评测集上的结果。
\end{itemize}


核心表可以这样设计：


{\footnotesize
\begin{tabular}{|p{\dimexpr 0.142\columnwidth-2\tabcolsep\relax}|p{\dimexpr 0.738\columnwidth-2\tabcolsep\relax}|p{\dimexpr 0.121\columnwidth-2\tabcolsep\relax}|}
\hline
\textbf{表名} & \textbf{关键字段} & \textbf{作用} \\
\hline
\texttt{ai\_prompt\_template} & \texttt{id}、\texttt{tenant\_id}、\texttt{name}、\texttt{scene\_code}、\texttt{type}、\texttt{status} & 管理 Prompt 逻辑名称 \\
\hline
\texttt{ai\_prompt\_version} & \texttt{id}、\texttt{template\_id}、\texttt{version\_no}、\texttt{content}、\texttt{variables\_schema}、\texttt{model\_config}、\texttt{created\_by}、\texttt{change\_reason} & 保存可回放的 Prompt 内容 \\
\hline
\texttt{ai\_prompt\_release} & \texttt{id}、\texttt{template\_id}、\texttt{version\_id}、\texttt{env}、\texttt{traffic\_ratio}、\texttt{tenant\_scope}、\texttt{status} & 控制灰度和回滚 \\
\hline
\texttt{ai\_prompt\_run} & \texttt{id}、\texttt{request\_id}、\texttt{version\_id}、\texttt{variables\_hash}、\texttt{input\_tokens}、\texttt{output\_tokens}、\texttt{created\_at} & 连接线上请求与 Prompt 版本 \\
\hline
\end{tabular}
}



变量注入时要避免两个坑：


\begin{enumerate}[itemsep=1pt, topsep=2pt]
  \item \textbf{变量未经清洗直接拼接}：用户输入、工具结果、检索片段都可能携带注入指令。应该用明确的分区标签和转义策略隔离。
  \item \textbf{Prompt 版本和代码版本脱节}：Prompt 里新增了变量，代码没传，线上直接生成空上下文。建议用 \texttt{variables\_schema} 做运行时校验。
\end{enumerate}


还有一个落库细节：\texttt{ai\_prompt\_run} 里通常只存变量摘要、Hash、Token 和关联 ID。完整用户输入、检索片段、工具返回如果包含 PII 或业务敏感信息，要按安全等级决定是否脱敏、加密、缩短留存周期，不能为了回放方便把所有明文都塞进表里。



一个最小接口示例：



\begin{lstlisting}[language=Java]
public interface PromptService {

    RenderedPrompt render(RenderPromptCommand command);

    PromptVersion publish(PublishPromptCommand command);

    void rollback(String templateId, String targetVersionId);
}
\end{lstlisting}



如果 Prompt 输出要被程序稳定解析，最好不要只靠“请返回 JSON”。结构化输出、JSON Schema、Function Calling 的工程细节可以看 大模型结构化输出：从 JSON 契约到 Function Calling 落地。



\section*{模型网关：多模型路由、fallback 与成本控制}


模型网关很容易被低估。很多团队一开始直接在业务代码里调用某个供应商 SDK，等到要换模型、做灰度、查成本时才发现处处耦合。



\subsection*{模型网关策略对比}

{\footnotesize
\begin{tabular}{|p{\dimexpr 0.154\columnwidth-2\tabcolsep\relax}|p{\dimexpr 0.308\columnwidth-2\tabcolsep\relax}|p{\dimexpr 0.292\columnwidth-2\tabcolsep\relax}|p{\dimexpr 0.246\columnwidth-2\tabcolsep\relax}|}
\hline
\textbf{策略} & \textbf{核心逻辑} & \textbf{适合场景} & \textbf{风险} \\
\hline
固定模型 & 某个场景固定调用一个模型 & 早期系统、低复杂度任务 & 成本和稳定性受单供应商影响 \\
\hline
成本优先路由 & 默认走低成本模型，失败或低置信度再升级 & 分类、摘要、轻量问答 & 低成本模型误判会传导到下游 \\
\hline
质量优先路由 & 高价值请求优先走高能力模型 & 法务、金融、医疗辅助、复杂 Agent & 成本高，需要预算控制 \\
\hline
延迟优先路由 & 按 P95/P99 延迟和可用区选择模型 & 实时聊天、语音、在线客服 & 可能牺牲复杂推理质量 \\
\hline
多模型投票 & 多模型并行生成，再由评审器选择 & 高风险内容、关键报告 & 成本和延迟都高 \\
\hline
fallback 链 & 主模型失败后切备用模型 & 大多数生产系统 & 备用模型能力差异会影响输出一致性 \\
\hline
\end{tabular}
}



\subsection*{Token 预算怎么做}


模型网关至少要在调用前做一次预算：



\begin{lstlisting}
预计输入 Token = System Prompt + 用户输入 + 历史消息 + RAG 片段 + Memory + Tool Schema
预计总 Token = 预计输入 Token + 最大输出 Token
\end{lstlisting}



如果超预算，别直接截断字符串。更稳的降级顺序是：


\begin{enumerate}[itemsep=1pt, topsep=2pt]
  \item 删除低相关 RAG 片段。
  \item 压缩早期历史消息。
  \item 减少工具 Schema，只保留候选工具。
  \item 降低最大输出长度。
  \item 切换长上下文模型。
  \item 拒绝执行并提示用户缩小范围。
\end{enumerate}


这里的 Token 预算和上下文压缩，和前面提到的 Context Engineering 是同一类问题。更完整的上下文装配、按需加载和降级策略，可以看 上下文工程（Context Engineering）是什么？和 Prompt Engineering 有什么区别？。



OpenTelemetry 文档里的 GenAI registry 能看到 \texttt{gen\_ai.request.model}、\texttt{gen\_ai.response.model}、\texttt{gen\_ai.usage.input\_tokens}、\texttt{gen\_ai.usage.output\_tokens}、\texttt{gen\_ai.response.time\_to\_first\_chunk}、retrieval、tool 等字段。不过 OpenTelemetry 站内也提示 GenAI 语义约定已迁移到独立仓库，落地时不要只复制旧文档里的字段名，应锁定当前版本并做字段映射。无论选择 Langfuse、LangSmith 还是自建观测平台，都建议尽量向通用字段靠拢，后续迁移和统一监控更轻松。



\section*{工具调用与权限：让模型只提出动作，系统决定能不能做}


Tool Calling 很容易让人产生错觉：模型返回了一个函数名和参数，系统执行就行。



这在生产环境很危险。



更稳的心智模型是：\textbf{模型只能提出“想调用什么工具”，真正执行前必须经过系统校验}。



工具运行时至少要包含 6 道关：


{\footnotesize
\begin{tabular}{|p{\dimexpr 0.108\columnwidth-2\tabcolsep\relax}|p{\dimexpr 0.892\columnwidth-2\tabcolsep\relax}|}
\hline
\textbf{环节} & \textbf{作用} \\
\hline
工具注册 & 声明工具名称、描述、参数 Schema、权限标签、风险等级 \\
\hline
工具检索 & 从大量工具中选出当前任务相关的少数工具，避免上下文膨胀 \\
\hline
参数校验 & 用 JSON Schema 或强类型对象校验必填、格式、枚举、范围 \\
\hline
权限校验 & 按用户、租户、角色、资源 ID 做后端鉴权 \\
\hline
二次确认 & 删除、支付、发送消息、改配置等敏感操作必须让用户确认 \\
\hline
审计日志 & 记录模型建议、最终参数、执行人、执行结果和回滚信息 \\
\hline
\end{tabular}
}



Anthropic、OpenAI 和 Google 的官方工具/函数调用文档都强调工具定义、参数结构和调用处理；Google 的文档还明确提醒，对会发送订单、更新数据库等有明显后果的函数调用，要在执行前让用户确认。工程实现增加硬规则：\textbf{模型不得代替系统做权限判断}。



即使供应商提供 server-side tool，业务侧也不能省掉自己的 ACL、审计和确认流。供应商负责把工具能力接进模型，业务系统负责判断这个用户、这个租户、这个资源在当前场景下能不能执行。



相关阅读：大模型结构化输出：从 JSON 契约到 Function Calling 落地；需要跨模型、Agent 或 IDE 复用工具时，见 什么是 Model Context Protocol（MCP）？和 Function Calling、Agent 什么关系？。



\begin{figure}[H]
\centering
\includegraphics[width=0.85\textwidth]{/home/mao/demos/latex_demos/ai-agent-knowledge/HTML/AI 应用系统设计：从 Prompt Demo 到生产级架构 _ JavaGuide_files/structured-output-function-calling-function-calling-pipeline.jpg}
\caption*{Function Calling 完整调用链路：模型只生成调用意图，真正执行工具的是业务侧}
\end{figure}



工具接口可以这样定义：



\begin{lstlisting}[language=Java]
public interface AiTool {

    ToolDefinition definition();

    ToolResult execute(ToolExecutionContext context, Map<String, Object> arguments);
}
\end{lstlisting}



工具定义里要有风险等级：



\begin{lstlisting}[language=Java]
public enum ToolRiskLevel {
    READ_ONLY,
    WRITE_LOW_RISK,
    WRITE_HIGH_RISK
}
\end{lstlisting}



对于 \texttt{WRITE\_HIGH\_RISK}，编排层必须把工具调用转换成“待确认动作”，不能直接执行。



\begin{figure}[H]
\centering
\includegraphics[width=0.85\textwidth]{/home/mao/demos/latex_demos/ai-agent-knowledge/HTML/AI 应用系统设计：从 Prompt Demo 到生产级架构 _ JavaGuide_files/structured-output-function-calling-tool-call-security.jpg}
\caption*{工具调用安全风险分层：按风险等级匹配不同的控制策略}
\end{figure}



\section*{RAG 与 Memory：共享知识和个性化记忆怎么协作}


RAG 和 Memory 都会把外部信息塞进上下文，但它们的治理方式不同。



\subsection*{一次请求里的协作顺序}


一次请求里的推荐顺序如下：


\begin{enumerate}[itemsep=1pt, topsep=2pt]
  \item 入口层确认用户身份和权限范围。
  \item Memory 服务在用户范围内检索偏好和长期事实。
  \item RAG 服务在租户和资源权限范围内检索共享知识库。
  \item Context 管理层对两类结果分别去重、过滤、压缩。
  \item 编排层把 Memory 放进“用户背景”区域，把 RAG 放进“证据资料”区域。
  \item 模型输出时要求区分“基于资料的事实”和“基于用户偏好的表达方式”。
\end{enumerate}


这套顺序主要是为了避免上下文污染。具体项目也可以先查 RAG 再查 Memory，但权限范围必须先确定，不能把“先检索、后过滤”当成默认方案。



\subsection*{怎么避免上下文污染}

{\footnotesize
\begin{tabular}{|p{\dimexpr 0.175\columnwidth-2\tabcolsep\relax}|p{\dimexpr 0.286\columnwidth-2\tabcolsep\relax}|p{\dimexpr 0.540\columnwidth-2\tabcolsep\relax}|}
\hline
\textbf{污染类型} & \textbf{典型表现} & \textbf{防护方式} \\
\hline
RAG 噪声污染 & 检索到无关文档，模型被带偏 & Hybrid Search、Rerank、Top-N 压缩、引用校验 \\
\hline
权限污染 & 用户拿到无权访问的文档片段 & 检索前 ACL 过滤，租户隔离，审计召回结果 \\
\hline
Memory 错误固化 & 用户一次临时说法被当成长期偏好 & 写入置信度、过期时间、用户可编辑、人工复核 \\
\hline
新旧事实冲突 & 旧版本制度和新版本制度同时进入上下文 & 版本字段、时间过滤、冲突检测 \\
\hline
Prompt 注入污染 & 文档里写着“忽略前面规则” & 文档内容分区、指令优先级、注入检测 \\
\hline
\end{tabular}
}



实践中，RAG 和 Memory 的结果不要直接拼成一段“背景资料”。要给模型清晰标注来源、时间、权限和可信度。模型看到的上下文越有结构，越不容易把“用户偏好”“公司制度”“工具结果”混成一类信息。



知识库不是一次导入就结束。文档版本、增量同步、去重、回滚和全量重建都会影响线上答案，具体可以看 RAG 知识库文档如何更新：增量更新、版本控制、去重与全量重建。如果问题需要跨文档关系、实体关系和全局摘要，传统向量检索不一定够用，可以继续看 万字详解 GraphRAG：为什么只靠向量检索撑不起复杂知识问答。



\section*{可观测与评测：没有回放，就没有优化}


\subsection*{Trace 应该记录什么}


AI 应用排查问题时，最怕只看到最终答案。



一次完整请求至少要记录这些数据：


{\footnotesize
\begin{tabular}{|p{\dimexpr 0.140\columnwidth-2\tabcolsep\relax}|p{\dimexpr 0.860\columnwidth-2\tabcolsep\relax}|}
\hline
\textbf{类别} & \textbf{建议记录} \\
\hline
Prompt & 模板名、版本、变量摘要、最终渲染后的消息结构 \\
\hline
检索 & Query、召回片段、分数、来源、权限过滤结果、Rerank 排名 \\
\hline
Memory & 命中的记忆、记忆来源、更新时间、置信度 \\
\hline
Tool & 工具名称、参数、权限结果、执行耗时、返回摘要、错误 \\
\hline
模型 & 供应商、模型名、采样参数、输入输出 Token、finish reason \\
\hline
延迟 & 入口耗时、检索耗时、模型 TTFT、总耗时、工具耗时 \\
\hline
成本 & 输入成本、输出成本、缓存命中、按租户和场景归因 \\
\hline
结果 & 最终答案、结构化解析结果、用户反馈、评测分数 \\
\hline
\end{tabular}
}



Langfuse、LangSmith、Google Vertex AI 和 OpenTelemetry 的官方文档都包含 tracing、datasets、evaluators、token usage、latency 等对象。工具可不同，需采集的信号大体相同。



\subsection*{评测应该怎么做}


评测体系可以单独成一条工程线。Golden Set 怎么建、RAG 和 Agent 指标怎么拆、LLM-as-Judge 怎么接入 CI，建议看 AI 应用评测体系：从 Golden Set 构建到线上灰度闭环。



评测别只问“答案好不好”。更可控的做法是拆成链路指标。不同平台对指标的命名不完全一样，下面这组更适合当内部指标口径：


\begin{itemize}[itemsep=1pt, topsep=2pt]
  \item \textbf{Context Recall}：正确证据有没有被召回。
  \item \textbf{Context Precision}：放进上下文的片段有多少是有用的。
  \item \textbf{Faithfulness}：答案是否忠于给定证据。
  \item \textbf{Answer Relevancy}：答案是否回应了用户问题。
  \item \textbf{Tool Success Rate}：工具调用是否成功完成。
  \item \textbf{Format Valid Rate}：结构化输出是否能被解析。
  \item \textbf{Cost per Success}：每次成功回答的平均成本。
\end{itemize}


LLM-as-Judge 可以用于自动评测，但不能当唯一裁判。它适合做大规模初筛、回归对比和线上抽样，关键业务仍要保留人工复核、规则校验和用户反馈。OpenAI、Google、Langfuse 这类平台的评测能力更新很快，甚至可能出现接口迁移或旧平台弃用，生产系统最好把“评测任务、评测样本、评测结果”沉淀在自己的数据模型里，外部平台作为执行器或看板接入。



一个实用闭环是：



\begin{lstlisting}
线上失败样本 -> 进入数据集 -> 固定版本回放 -> 定位 Prompt/RAG/Tool/模型问题 -> 灰度新策略 -> 对比指标 -> 再发布
\end{lstlisting}



没有回放，就只能靠感觉调 Prompt。靠感觉调出来的系统，线上很难稳住。



\section*{安全与合规：AI 应用的风险入口更多}


AI 应用的安全面比传统 CRUD 系统更宽。因为用户输入、检索文档、工具返回、历史记忆都可能影响模型行为。



\subsection*{风险项要落到代码和流程里}

{\footnotesize
\begin{tabular}{|p{\dimexpr 0.161\columnwidth-2\tabcolsep\relax}|p{\dimexpr 0.482\columnwidth-2\tabcolsep\relax}|p{\dimexpr 0.357\columnwidth-2\tabcolsep\relax}|}
\hline
\textbf{风险} & \textbf{说明} & \textbf{处理建议} \\
\hline
PII 泄露 & 日志、Prompt、评测集里包含手机号、身份证、邮箱等 & 入库前脱敏，敏感字段加密，最小化留存 \\
\hline
权限绕过 & 检索或工具调用绕过业务 ACL & 检索前过滤，工具执行前二次鉴权 \\
\hline
Prompt 注入 & 用户或文档诱导模型忽略系统规则 & 内容分区、指令优先级、注入检测、拒答策略 \\
\hline
数据留存失控 & 模型请求和观测日志保存过久 & 按租户和场景配置留存周期 \\
\hline
训练数据风险 & 把用户敏感数据用于微调或评测 & 明确授权、脱敏、隔离、可删除 \\
\hline
高风险动作误执行 & 模型误调用删除、支付、发信等工具 & 风险分级、二次确认、审计和补偿 \\
\hline
\end{tabular}
}



这里有个容易忽略的细节：\textbf{安全策略不能只写在 Prompt 里}。Prompt 可以提醒模型“不要泄露隐私”，但权限过滤、脱敏、审计、确认流必须由代码和基础设施强制执行。



\subsection*{第三方模型要单独管数据边界}


如果请求会发往第三方模型，还要单独确认数据授权、区域、留存和训练使用策略。拿不准时，默认按最小化原则处理：能不发的字段不发，必须发的字段先脱敏或摘要化，并把留存周期写进配置和审计里。



Prompt 注入、上下文分区和工具权限是连在一起的，前面提到的 Prompt Engineering、Context Engineering 和 MCP 这几篇可以配合看。



\section*{Java 后端落地建议}


用 Java 构建生产级 AI 应用时，按“领域能力”拆模块，不要按供应商 SDK 拆模块。



\subsection*{模块拆分}

{\footnotesize
\begin{tabular}{|p{\dimexpr 0.346\columnwidth-2\tabcolsep\relax}|p{\dimexpr 0.654\columnwidth-2\tabcolsep\relax}|}
\hline
\textbf{模块} & \textbf{职责} \\
\hline
\texttt{ai-api} & 对外 REST/SSE/WebSocket 接口，请求鉴权和协议适配 \\
\hline
\texttt{ai-orchestrator} & 业务编排、交互模式选择、任务状态机 \\
\hline
\texttt{ai-prompt} & Prompt 模板、版本、灰度、渲染、回滚 \\
\hline
\texttt{ai-context} & 上下文组装、Token 预算、历史压缩、上下文分区 \\
\hline
\texttt{ai-gateway} & 模型路由、fallback、限流、熔断、成本统计 \\
\hline
\texttt{ai-rag} & 知识库检索、权限过滤、Rerank、引用管理 \\
\hline
\texttt{ai-memory} & 用户记忆写入、检索、冲突处理、过期策略 \\
\hline
\texttt{ai-tool} & 工具注册、参数校验、执行、二次确认、审计 \\
\hline
\texttt{ai-eval} & 数据集、评测任务、LLM-as-Judge、人工反馈 \\
\hline
\texttt{ai-observability} & Trace、指标、日志、成本、告警 \\
\hline
\end{tabular}
}



\subsection*{核心表设计}


这组表不要求第一版全部建完，它主要说明生产系统里哪些数据要有归属。第一版至少要把请求 Trace、模型调用、Prompt 版本、RAG 召回记录落下来，后面排查问题才有材料。


{\footnotesize
\begin{tabular}{|p{\dimexpr 0.092\columnwidth-2\tabcolsep\relax}|p{\dimexpr 0.738\columnwidth-2\tabcolsep\relax}|p{\dimexpr 0.169\columnwidth-2\tabcolsep\relax}|}
\hline
\textbf{表名} & \textbf{建议关键字段} & \textbf{作用} \\
\hline
\texttt{ai\_request\_trace} & \texttt{id}、\texttt{request\_id}、\texttt{tenant\_id}、\texttt{user\_id}、\texttt{scene\_code}、\texttt{mode}、\texttt{status}、\texttt{total\_latency\_ms}、\texttt{error\_code}、\texttt{created\_at} & 一次 AI 请求的主 Trace，记录用户、租户、场景、状态、耗时 \\
\hline
\texttt{ai\_model\_call} & \texttt{id}、\texttt{request\_id}、\texttt{provider}、\texttt{model\_name}、\texttt{prompt\_version\_id}、\texttt{input\_tokens}、\texttt{output\_tokens}、\texttt{ttft\_ms}、\texttt{latency\_ms}、\texttt{finish\_reason}、\texttt{error\_code} & 模型调用明细，记录模型、参数、Token、TTFT、错误 \\
\hline
\texttt{ai\_context\_item} & \texttt{id}、\texttt{request\_id}、\texttt{source\_type}、\texttt{source\_id}、\texttt{content\_hash}、\texttt{token\_count}、\texttt{inject\_position}、\texttt{sensitivity\_level} & 上下文条目，记录来源类型、来源 ID、Token、注入位置 \\
\hline
\texttt{ai\_rag\_chunk\_hit} & \texttt{id}、\texttt{request\_id}、\texttt{knowledge\_base\_id}、\texttt{doc\_id}、\texttt{chunk\_id}、\texttt{score}、\texttt{rank\_no}、\texttt{acl\_result}、\texttt{citation\_url} & RAG 召回明细，记录分数、排名、文档权限、引用信息 \\
\hline
\texttt{ai\_memory\_item} & \texttt{id}、\texttt{tenant\_id}、\texttt{user\_id}、\texttt{memory\_type}、\texttt{content}、\texttt{confidence}、\texttt{expires\_at}、\texttt{status}、\texttt{updated\_at} & 长期记忆条目，记录用户、内容、置信度、过期时间、状态 \\
\hline
\texttt{ai\_tool\_call} & \texttt{id}、\texttt{request\_id}、\texttt{tool\_name}、\texttt{risk\_level}、\texttt{arguments\_hash}、\texttt{permission\_result}、\texttt{confirm\_status}、\texttt{execute\_status}、\texttt{latency\_ms} & 工具调用明细，记录工具、参数摘要、权限结果、执行结果 \\
\hline
\texttt{ai\_eval\_dataset} & \texttt{id}、\texttt{name}、\texttt{scene\_code}、\texttt{version\_no}、\texttt{status}、\texttt{created\_by} & 评测集元信息 \\
\hline
\texttt{ai\_eval\_case} & \texttt{id}、\texttt{dataset\_id}、\texttt{input}、\texttt{expected\_behavior}、\texttt{tags}、\texttt{difficulty}、\texttt{status} & 评测样本，包含输入、期望行为、标签 \\
\hline
\texttt{ai\_eval\_run} & \texttt{id}、\texttt{dataset\_id}、\texttt{target\_type}、\texttt{target\_version}、\texttt{judge\_config}、\texttt{status}、\texttt{started\_at}、\texttt{finished\_at} & 某次评测任务 \\
\hline
\texttt{ai\_eval\_result} & \texttt{id}、\texttt{run\_id}、\texttt{case\_id}、\texttt{score}、\texttt{pass\_status}、\texttt{judge\_reason}、\texttt{error\_code} & 单条样本评测结果 \\
\hline
\end{tabular}
}



表设计里有 3 个细节别省：


\begin{enumerate}[itemsep=1pt, topsep=2pt]
  \item \texttt{request\_id} 要贯穿 Prompt、RAG、Memory、Tool、Model Call 和 Eval，最好全链路唯一。
  \item 大字段不要无脑进 MySQL。完整 Prompt、模型输出、工具返回可以放对象存储或日志系统，业务表里保留摘要、Hash、敏感级别和引用地址。
  \item 运行时表要按 \texttt{tenant\_id}、\texttt{scene\_code}、\texttt{created\_at}、\texttt{status} 设计索引和归档策略，否则观测表很快会变成新的性能瓶颈。
\end{enumerate}


\subsection*{核心接口设计}


\begin{lstlisting}[language=Java]
public interface ModelGateway {

    ModelResponse generate(ModelRequest request);

    Flux<ModelStreamEvent> stream(ModelRequest request);
}
\end{lstlisting}



如果项目没有用 WebFlux，\texttt{Flux<ModelStreamEvent>} 可以替换成 JDK \texttt{Flow.Publisher}、SSE emitter 或内部事件回调。重点是把“同步生成”和“流式事件”分成两个接口语义，不要让调用方猜返回值到底什么时候完整。



\begin{lstlisting}[language=Java]
public interface ContextAssembler {

    AssembledContext assemble(AiRequest request, ContextPolicy policy);
}
\end{lstlisting}



\begin{lstlisting}[language=Java]
public interface RagService {

    List<RagHit> retrieve(RagQuery query, PermissionScope permissionScope);
}
\end{lstlisting}



\begin{lstlisting}[language=Java]
public interface EvaluationService {

    EvalRunResult runDataset(EvalRunCommand command);
}
\end{lstlisting}



\subsection*{一个最小请求链路}


\begin{lstlisting}
Controller
  -> RequestGuard 鉴权、限流、脱敏
  -> Orchestrator 选择同步/流式/异步
  -> ContextAssembler 拉取 RAG、Memory、历史
  -> PromptService 渲染模板版本
  -> ModelGateway 路由模型并记录 Token
  -> OutputParser 校验结构化输出
  -> TraceService 写入观测数据
\end{lstlisting}



仅构建企业知识库问答时，第一阶段可先落地 \texttt{ai-api}、\texttt{ai-prompt}、\texttt{ai-gateway}、\texttt{ai-rag}、\texttt{ai-observability}。Memory、Tool、Eval 可以逐步补齐。但 Trace 和 Prompt 版本不要拖到后面，它们是后续排查问题的地基。



需要从 Java 后端调用大模型 API 的细节入手时，先看 大模型 API 调用工程实践；团队准备把模型调用统一成基础设施时，单独读一遍 大模型网关详解。



\section*{要点回顾}

\begin{enumerate}[itemsep=1pt, topsep=2pt]
  \item \textbf{Prompt Demo 只证明“能回答”，生产级架构要证明“长期可控地回答”}。
  \item \textbf{模型网关是 AI 应用的模型调用控制面}，负责路由、fallback、限流、熔断、Token 预算和成本归因。
  \item \textbf{Prompt 必须版本化}，支持变量校验、灰度、回滚和审计。
  \item \textbf{RAG、Memory、Tool 要分开治理}，共享知识、个性化记忆和真实业务动作不能混成一团。
  \item \textbf{可观测和评测决定系统能不能持续变好}，没有 Trace 和回放，优化基本靠猜。
  \item \textbf{安全策略要靠代码强制执行}，Prompt 只能辅助，不能替代权限、脱敏、审计和二次确认。
\end{enumerate}


\section*{参考资料}


JavaGuide 相关阅读：


\begin{itemize}[itemsep=1pt, topsep=2pt]
  \item AI 应用开发知识体系：大模型、Agent、RAG、MCP、Prompt 工程与系统设计
  \item AI 系统设计专题：生产级架构、模型网关、评测治理与语音 Agent
  \item 大模型基础专题：运行机制、API 调用、结构化输出与评测
  \item RAG 专题：文档处理、向量数据库、GraphRAG、检索优化与知识库更新
  \item AI Agent 专题：Agent Loop、Memory、Prompt、Context、MCP 与 Skills
  \item \href{https://developers.openai.com/api/docs}{OpenAI API 官方文档}
  \item \href{https://developers.openai.com/api/docs/guides/function-calling}{OpenAI Function Calling 官方文档}
  \item \href{https://developers.openai.com/api/docs/guides/streaming-responses}{OpenAI Streaming 官方文档}
  \item \href{https://developers.openai.com/api/docs/guides/evals}{OpenAI Evals 官方文档}
  \item \href{https://developers.openai.com/api/docs/guides/agents/integrations-observability}{OpenAI Agents SDK 观测与集成}
  \item \href{https://platform.claude.com/docs/en/agents-and-tools/tool-use/overview}{Anthropic Tool Use 官方文档}
  \item \href{https://platform.claude.com/docs/en/build-with-claude/prompt-caching}{Anthropic Prompt Caching 官方文档}
  \item \href{https://docs.cloud.google.com/gemini-enterprise-agent-platform/models/tools/function-calling}{Google Gemini Function Calling 官方文档}
  \item \href{https://docs.cloud.google.com/gemini-enterprise-agent-platform/models/evaluation-overview}{Google 生成式 AI 评测官方文档}
  \item \href{https://docs.cloud.google.com/gemini-enterprise-agent-platform/models/grounding/ground-responses-using-rag}{Google RAG Grounding 官方文档}
  \item \href{https://langfuse.com/docs/observability/overview}{Langfuse Observability 官方文档}
  \item \href{https://langfuse.com/docs/prompt-management/overview}{Langfuse Prompt Management 官方文档}
  \item \href{https://docs.langchain.com/langsmith/evaluation}{LangSmith Evaluation 官方文档}
  \item \href{https://opentelemetry.io/docs/specs/semconv/registry/attributes/gen-ai/}{OpenTelemetry GenAI 属性注册表}
\end{itemize}

\end{document}
